\documentclass[11pt]{article}

\usepackage[margin=1in]{geometry}
\usepackage{amsmath,amssymb}
\usepackage{array}
\usepackage{booktabs}
\usepackage{hyperref}
\usepackage{xcolor}
\usepackage{enumitem}
\usepackage{listings}

\hypersetup{
  colorlinks=true,
  linkcolor=blue!60!black,
  urlcolor=blue!60!black,
  citecolor=blue!60!black
}

\newcommand{\code}[1]{\texttt{\detokenize{#1}}}
\newcommand{\R}{\mathbb{R}}
\newcommand{\E}{\mathcal{E}}
\newcommand{\K}{\mathcal{K}}
\newcommand{\G}{\mathcal{G}}
\newcommand{\V}{\mathcal{V}}
\newcommand{\A}{\mathcal{A}}

\lstset{
  basicstyle=\ttfamily\small,
  columns=fullflexible,
  keepspaces=true,
  frame=single,
  breaklines=true,
  backgroundcolor=\color{black!3}
}

\title{Upwind SCC Graph Ordering for HDG Trace Systems}
\author{\code{hdgfem.linalg.ordering}}
\date{\today}

\begin{document}
\maketitle

\begin{abstract}
This note documents the ordering algorithm implemented in
\code{linalg/ordering.py}.  The purpose of the module is to build an
experimental permutation for HDG trace unknowns in advection-dominated systems.
The ordering is algebraic with respect to the sparse trace system, but it is
constructed from mesh-edge connectivity and the sign of the normal advective
flux on element faces.  The main pipeline is:
\begin{enumerate}[nosep]
  \item build a directed graph between mesh-edge blocks using local
        inflow-to-outflow dependencies;
  \item compute strongly connected components (SCCs) with Kosaraju's
        algorithm;
  \item topologically order the SCC condensation graph;
  \item lift the ordered edge blocks to a trace degree-of-freedom permutation.
\end{enumerate}
The module also contains sparsity-pattern plotting utilities for comparing the
matrix before and after the upwind SCC permutation.
\end{abstract}

\tableofcontents

\section{Scope and Design Intent}

The HDG global unknown is the trace field on mesh edges.  For a uniform trace
polynomial degree, each global mesh edge carries the same number of trace
degrees of freedom, denoted here by
\[
  n_{\partial} = \code{edg_dof}.
\]
The graph-ordering module does not alter element assembly, mesh connectivity,
or numerical fluxes.  It only computes a permutation of trace unknown indices.
The convention matches \code{hdgfem.linalg.system.solve_global_system}:
\[
  A_{\pi} = A[\pi, \pi], \qquad b_{\pi} = b[\pi],
\]
and after solving the permuted system the solution is unpermuted by
\[
  x[\pi] = x_{\pi}.
\]

The ordering is designed to expose approximate upstream-to-downstream
structure in advection-dominated HDG trace matrices.  It should be viewed as a
preconditioner-ordering experiment, not as a guaranteed fill-reducing ordering.
In particular, SuperLU's internal orderings such as \code{COLAMD} may still be
better at high polynomial order because they optimize symbolic fill rather than
flow direction.

\section{Inputs and Core Data Structures}

\subsection{Mesh edge connectivity}

The primary mesh input is the element-local to global-edge map
\[
  \code{loc2glob_edge}[K,f] \in \{0,\ldots,N_e-1\},
\]
where \(K\) is an element index, \(f\in\{0,1,2\}\) is a local triangle face, and
\(N_e\) is the number of global mesh edges.  In the code this is passed to the
Numba kernel as \code{loc2glob_edge}.  For boundary-eliminated trace systems,
only non-boundary edge blocks remain active.  The helper \code{_as_active_edges}
builds three arrays:
\begin{itemize}[nosep]
  \item \code{active_edge_ids}: sorted global edge ids participating in the
        reduced system;
  \item \code{active_mask}: Boolean mask over all mesh edges;
  \item \code{old_to_active}: map from global edge id to active graph node id,
        with \(-1\) for inactive edges.
\end{itemize}

\subsection{Face-normal advective flux}

The second geometric/numerical input is
\[
  \code{beta_dot_normal}[K,f,q]
  \approx \beta_h(x_{K,f,q})\cdot n_{K,f},
\]
where \(q\) indexes face quadrature points.  The graph construction reduces
this to one mean face flux per element-local face:
\[
  \bar F_{K,f}
  =
  \frac{1}{N_q}\sum_{q=1}^{N_q}
  \beta_h(x_{K,f,q})\cdot n_{K,f}.
\]
A face is classified as inflow when
\[
  \bar F_{K,f} < -\varepsilon,
\]
and as outflow when
\[
  \bar F_{K,f} > \varepsilon,
\]
where \(\varepsilon=\code{flux_tolerance}\).  Faces in the neutral band
\([-\varepsilon,\varepsilon]\) are ignored.

\subsection{Return containers}

The algorithm returns a \code{GraphOrderingResult}:
\[
\begin{array}{ll}
\code{edge_order} & \text{active global edge ids in graph order},\\
\code{dof_permutation} & \text{trace dof permutation used by the sparse solver},\\
\code{component_id} & \text{SCC id for each active graph node},\\
\code{component_order} & \text{topological order of SCC ids},\\
\code{diagnostics} & \text{graph sizes and timing information}.
\end{array}
\]
The diagnostics contain the number of graph nodes, directed graph edges, SCCs,
the largest SCC size, the number of cyclic SCCs, the number of nodes in cyclic
SCCs, and a timing breakdown.

\section{Step 1: Build the Directed Upwind Edge Graph}

\subsection{Graph nodes}

The directed graph nodes are active mesh-edge blocks, not individual trace
degrees of freedom.  If boundary dofs have not been eliminated, every mesh edge
is active.  If boundary dofs have been eliminated, only non-boundary edges are
active.  Therefore the graph vertex set is
\[
  \V = \{0,\ldots,N_a-1\},
\]
where \(N_a=\code{active_edge_ids.size}\).

\subsection{Element-local dependencies}

For each element \(K\), the kernel inspects all local faces.  Every active
inflow face is connected to every active outflow face in the same element:
\[
  e_{K,f_{\rm in}} \longrightarrow e_{K,f_{\rm out}}.
\]
In code, this is done by \code{_build_upwind_edge_pairs_kernel}.  In words:

\begin{lstlisting}
for each element K:
    compute mean face flux F[f] for f=0,1,2
    for each source face f:
        keep it only if active and F[f] < -flux_tolerance
        for each target face g != f:
            keep it only if active and F[g] > flux_tolerance
            append edge active_id(f) -> active_id(g)
\end{lstlisting}

The direction ``inflow to outflow'' is the local transport dependency direction:
information entering an element through an inflow face contributes to unknowns
on outflow faces.  The graph can contain duplicate directed edges.  Duplicates
are accepted because they do not affect reachability or SCC membership, and
removing them would add a sorting/unique phase.

\subsection{Storage}

The kernel emits two arrays
\[
  \code{sources},\quad \code{targets}
\]
with equal length.  Directed graph edge \(m\) is
\[
  \code{sources}[m]\rightarrow \code{targets}[m].
\]
The maximum number of appended pairs is bounded by \(6N_K\), because each
triangle has three faces and at most each inflow face connects to the two other
faces.  Hence the temporary arrays are allocated with length
\[
  6\,\code{num_elements}.
\]

\section{Step 2: Build CSR Adjacency}

The helper \code{_build_csr_kernel} converts edge arrays into compressed sparse
row adjacency:
\[
  \code{indptr}\in\mathbb{N}^{N_a+1}, \qquad
  \code{indices}\in\mathbb{N}^{N_{\rm graph\ edges}}.
\]
The implementation has two passes:
\begin{enumerate}[nosep]
  \item count outgoing edges per source node and prefix-sum the counts;
  \item fill \code{indices} using a mutable copy of \code{indptr}.
\end{enumerate}
The same routine is used for the reverse graph by swapping sources and
targets.  Duplicate edges remain in the adjacency.

\section{Step 3: Strongly Connected Components}

\subsection{Why SCCs?}

If the upwind graph is acyclic, a topological ordering gives an exact upstream
to downstream order.  If the flow recirculates, some edge blocks are mutually
reachable.  These cycles cannot be linearly ordered without breaking a
dependency.  The algorithm therefore groups mutually reachable nodes into SCCs,
then orders the SCCs.  Inside one SCC, nodes are kept in their original order.

\subsection{Kosaraju algorithm}

The function \code{_kosaraju_scc_kernel} implements Kosaraju's two-pass SCC
algorithm using explicit stacks so that it works inside Numba without Python
recursion.

The first pass runs depth-first search on the original graph and records nodes
by finishing time:
\[
  \code{finish_order}[0],\ldots,\code{finish_order}[N_a-1].
\]
The second pass scans nodes in reverse finishing order and runs DFS on the
reverse graph.  Each reverse DFS discovers one SCC.  The outputs are:
\[
  \code{component_id}[v] = c,
  \qquad
  \code{component_sizes}[c] = |\{v:\code{component_id}[v]=c\}|.
\]

\subsection{Correctness property}

For any directed graph, Kosaraju's algorithm partitions vertices into maximal
sets of mutually reachable vertices.  Collapsing each SCC into one component
produces a directed acyclic graph.  That acyclicity is what makes the next
topological-order step well-defined.

\section{Step 4: Build and Order the SCC Condensation DAG}

\subsection{Component graph}

The function \code{_component_edges_kernel} maps each original graph edge
\[
  u\rightarrow v
\]
to an SCC edge
\[
  C(u)\rightarrow C(v).
\]
Self-edges, where \(C(u)=C(v)\), are discarded.  The result is the condensation
graph.  It is a DAG by construction.

\subsection{Topological order}

The function \code{_topological_order_kernel} computes indegrees in the
condensation DAG and applies Kahn's queue algorithm:
\begin{lstlisting}
queue = all components with indegree zero
while queue is not empty:
    pop a component c
    append c to the order
    for each outgoing edge c -> d:
        decrement indegree[d]
        if indegree[d] becomes zero:
            push d
\end{lstlisting}
The code checks that the number of ordered components equals the number of
components.  Failure would indicate a bug, because the condensation graph of an
SCC partition must be acyclic.

\section{Step 5: Lift Component Order to Edge-Block Order}

The function \code{_node_order_from_components_kernel} produces an ordered list
of graph nodes.  It reserves contiguous ranges for components according to the
topological component order, then inserts original node ids into the range of
their component.  Thus:
\begin{itemize}[nosep]
  \item components are topologically ordered;
  \item nodes inside one SCC keep original active-node order.
\end{itemize}
The active-node order is then converted back to global mesh-edge ids:
\[
  \code{edge_order}
  =
  \code{active_edge_ids}[\code{node_order}].
\]

\section{Step 6: Lift Edge Blocks to Trace DOFs}

The sparse trace matrix is indexed by individual trace degrees of freedom.  If
edge block \(e\) has \(n_{\partial}\) local trace dofs, its trace dof range is
\[
  e n_{\partial},\ e n_{\partial}+1,\ \ldots,\ e n_{\partial}+n_{\partial}-1.
\]
The helper \code{_trace_dof_permutation_kernel} receives ordered active-node
positions, not global edge ids.  For each ordered active node position \(a\),
it appends the contiguous local block
\[
  a n_{\partial},\ a n_{\partial}+1,\ldots,a n_{\partial}+n_{\partial}-1.
\]
This is exactly the correct indexing for reduced systems where active edge
blocks have been compacted into consecutive reduced-system coordinates.

For a full, non-reduced trace system, active nodes are all edges in natural
global edge order.  Then the active-node position and global edge id coincide.
For a boundary-eliminated system, \code{active_edges} must be passed so the
active edge positions match the reduced-system matrix.

\section{Boundary Elimination and Active Edges}

The ordering routine is used in two common modes:

\paragraph{Full penalty system.}
Boundary trace dofs remain in the system.  The active edge set is all mesh
edges:
\[
  \code{active_edges} = \code{None}.
\]

\paragraph{Boundary-eliminated system.}
Prescribed boundary trace dofs are removed before the solve.  The reduced
matrix contains only non-boundary edge blocks.  In this case the caller builds
\[
  \code{active_edges} = \{e:\ e \notin \partial\Omega\}.
\]
The graph is built using global edge ids, but all graph nodes are compacted
through \code{old_to_active}.  The resulting dof permutation is therefore in
reduced-system coordinates and can be applied directly to the reduced matrix.

\section{Complexity}

Let \(N_K\) be the number of elements, \(N_a\) the number of active edge
blocks, and \(M_g\) the number of directed graph edges after local inflow to
outflow construction.  For triangular elements,
\[
  M_g \le 6N_K.
\]
Ignoring duplicate graph edges, the algorithmic costs are:
\[
\begin{array}{ll}
\text{graph pair construction} & O(N_K N_q),\\
\text{CSR construction} & O(N_a + M_g),\\
\text{Kosaraju SCC} & O(N_a + M_g),\\
\text{component DAG construction} & O(M_g),\\
\text{topological order} & O(N_{\rm comp}+M_{\rm comp}),\\
\text{dof lift} & O(N_a n_{\partial}).
\end{array}
\]
The graph-ordering cost is usually small compared with HDG matrix assembly and
incomplete factorization.  It is also almost independent of the polynomial
degree \(p\), except for the final dof lift through \(n_{\partial}=p+1\).

\section{Diagnostics}

\code{GraphOrderingDiagnostics} reports:
\[
\begin{array}{ll}
\code{num_nodes} & N_a,\\
\code{num_directed_edges} & M_g,\\
\code{num_components} & N_{\rm comp},\\
\code{largest_component_size} & \max_c |C_c|,\\
\code{cyclic_components} & \#\{c:\ |C_c|>1\},\\
\code{cyclic_nodes} & \sum_{c:\ |C_c|>1}|C_c|.
\end{array}
\]
For a nearly acyclic upwind graph, one typically sees
\[
  \code{largest_component_size}=1,\qquad \code{cyclic_nodes}=0.
\]
For recirculating flows, SCCs larger than one indicate trace-edge blocks whose
flow dependencies are mutually coupled.

The timing object \code{GraphOrderingTimings} records the wall-clock time for
graph pair construction, CSR construction, SCC, DAG construction, topological
ordering, dof permutation, and total ordering.

\section{Sparse Pattern Plotting}

The module also provides two plotting helpers:
\[
\begin{array}{ll}
\code{save_sparse_pattern_plot} & \text{plot one sparse matrix pattern},\\
\code{save_upwind_reordered_matrix_patterns}
  & \text{plot before and after } A[\pi,\pi].
\end{array}
\]
The marker area is chosen by \code{sparse_pattern_marker_area}.  The intended
law is inverse in the number of nonzeros:
\[
  s
  =
  s_{\rm ref}\frac{N_{\rm ref}}{\operatorname{nnz}(A)}.
\]
The current implementation clips this value and applies an additional practical
floor, so very large matrices do not disappear completely in PNG output.  Very
large matrices are also stride-sampled to \code{max_plot_points}; the marker
area is still computed from the original \(\operatorname{nnz}(A)\), not from
the sampled count.

\section{Interaction with SuperLU ILU}

The upwind SCC permutation is not a symbolic fill-reducing ordering.  It is a
physics-motivated block ordering.  If the sparse solver is SuperLU's
\code{spilu}, then its \code{permc_spec} option matters:
\begin{itemize}
  \item \code{NATURAL} forces SuperLU to respect the supplied matrix ordering.
        This can expose the benefit of a good upwind order, but it can also be
        fragile at high polynomial degree.
  \item \code{COLAMD} lets SuperLU apply a fill-reducing column ordering.  This
        may override much of the upwind-ordering effect, but it can be far more
        robust when local trace blocks become large.
\end{itemize}
Thus the SCC ordering should be benchmarked together with the ILU column
permutation strategy.  A small edge-block graph can look excellent while the
lifted high-order dof matrix still has poor symbolic fill under \code{NATURAL}.

\section{High-Level Pseudocode}

\begin{lstlisting}
function upwind_scc_trace_ordering(mesh, beta_dot_normal, edg_dof,
                                   active_edges=None, flux_tolerance=0):
    active_edge_ids, active_mask, old_to_active = active_edges_to_maps(...)

    sources, targets =
        build_upwind_edge_pairs(mesh.loc2glob_edge,
                                beta_dot_normal,
                                active_mask,
                                old_to_active,
                                flux_tolerance)

    node_order, component_id, component_order, component_sizes, timings =
        strongly_connected_component_order(num_active_edges, sources, targets)

    edge_order = active_edge_ids[node_order]
    dof_permutation = lift_edge_order_to_trace_dofs(node_order, edg_dof)

    return GraphOrderingResult(edge_order,
                               dof_permutation,
                               component_id,
                               component_order,
                               diagnostics)
\end{lstlisting}

\section{Function Map}

\begin{center}
\small
\begin{tabular}{p{0.42\linewidth}p{0.48\linewidth}}
\toprule
Function & Role \\
\midrule
\code{_build_upwind_edge_pairs_kernel} & Build directed inflow-to-outflow edge graph. \\
\code{_build_csr_kernel} & Convert directed edge arrays to CSR adjacency. \\
\code{_kosaraju_scc_kernel} & Compute SCCs with iterative Kosaraju DFS. \\
\code{_component_edges_kernel} & Build condensation graph edges between SCCs. \\
\code{_topological_order_kernel} & Topologically order the condensation DAG. \\
\code{_node_order_from_components_kernel} & Convert component order to node order. \\
\code{_trace_dof_permutation_kernel} & Lift ordered edge blocks to trace dofs. \\
\code{strongly_connected_component_order} & Public graph-only SCC ordering routine. \\
\code{upwind_scc_trace_ordering} & Public HDG trace-edge ordering routine. \\
\code{sparse_pattern_marker_area} & Compute nnz-dependent marker area for plots. \\
\code{save_sparse_pattern_plot} & Save one matrix sparsity pattern figure. \\
\code{save_upwind_reordered_matrix_patterns} & Save before/after upwind permutation figures. \\
\bottomrule
\end{tabular}
\end{center}

\section{Practical Usage}

From the advection-reaction CLI, the current diagnostic path can generate the
matrix patterns without entering the expensive global solve:

\begin{lstlisting}
python -m hdgfem.solvers.adv_rea -p 7 --lc 0.03 \
  --boundary-mode eliminate \
  --matrix-pattern-only \
  --matrix-pattern-max-points 2000000
\end{lstlisting}

This writes figures such as:
\begin{lstlisting}
matrix_patterns/*_before_upwind.png
matrix_patterns/*_after_upwind.png
\end{lstlisting}

\section{Summary}

The ordering in \code{graph_ordering.py} is a block-level flow ordering:
construct local inflow-to-outflow dependencies, compress cycles with SCCs,
topologically order the acyclic component graph, and lift the result to trace
dofs.  The graph construction is inexpensive and mostly independent of
polynomial order.  Its effect on ILU depends on how the block ordering interacts
with the symbolic fill behavior of the sparse factorization.

\end{document}
